\documentclass[10pt,a4paper]{amsart}
\usepackage[margin=19mm]{geometry}
\usepackage{amsmath,amssymb,mathtools}
\usepackage[scr=rsfs,cal=euler]{mathalfa}
\usepackage{xfrac}
\usepackage[unicode,hidelinks,bookmarks=false]{hyperref}
\newcommand{\ie}{i.e.\ }
\newcommand{\cf}{cf.\ }
\newcommand{\resp}{respectively\ }
\newcommand{\Subsection}{\subsection}
\providecommand{\nobreakdash}{\nobreak\hbox{-}}
\setlength{\emergencystretch}{2em}
\multlinegap0pt
\setlength{\voffset}{-.25in}
\usepackage{latexsym}
\usepackage{color}
\usepackage{url}
\numberwithin{equation}{section}
\newtheorem{thm}{Theorem}[section]
\newtheorem{theorem}[thm]{Theorem}
\newtheorem{lemma}[thm]{Lemma}
\newtheorem{example}[thm]{Example}
\newtheorem{definition}[thm]{Definition}
\newtheorem{proposition}[thm]{Proposition}
\newtheorem{corollary}[thm]{Corollary}
\newtheorem{remark}[thm]{Remark}
\newtheorem{consequence}[thm]{Consequence}
\DeclareMathOperator{\arctanh}{arctanh}
\DeclareMathOperator{\arccot}{arccot}
\DeclareMathOperator{\Li}{Li}
\begin{document}
\title[A combinatorial sum with two complex parameters]{A combinatorial sum with two complex parameters}
\author{MSC 2000 : 05A10, 11B65, 11B83; Michel Bataille; Robert Frontczak}
\date{}
\maketitle
\noindent\emph{Source and licence.} A combinatorial sum with two complex parameters. Version arXiv:2607.02639v1 (CC BY 4.0). This material is distributed under the Creative Commons Attribution 4.0 International licence, http://creativecommons.org/licenses/by/4.0/, as recorded in the arXiv OAI licence metadata for this exact version and shown on the abstract page https://arxiv.org/abs/2607.02639v1. Full rights notice: licenses/arxiv-2607.02639-CC-BY-4.0.txt.\par\smallskip
\noindent\textit{Editorial scope.} Counted unit: the original \texttt{corollary} environment \texttt{from\_Cx1} at source lines 377--397 together with its complete proof at lines 398--400; the delivered document keeps source lines 83--524 so that every referenced label and every citation resolves inside it. Native editable source is the paper's own concatenated TeX; the AMS wrapper is editorial formatting only. No publisher class, upstream script or unrelated artwork is redistributed.
\begin{equation}\label{Gould_id}
f(x+y) = y \binom{y+n}{n} \sum_{k=0}^n \binom{n}{k} (-1)^k \frac{f(x-k)}{y+k}.
\end{equation}
If we set $y=x$ and $f(x)=1$ then
we obtain the well-known decomposition
\begin{equation}\label{partial}
\sum_{k=0}^n \binom{n}{k} \frac{(-1)^k}{x+k} = \frac{1}{x\binom{x+n}{n}} \qquad x\notin \{-n,-(n-1),\ldots,0\}.
\end{equation}
In this paper, we work with an identity similar in type to \eqref{Gould_id}. It involves two complex parameters and appears very productive. This is first shown by several applications leading to a number of identities and evaluations of combinatorial sums (Catalan's numbers intervene in some of them). Then the focus turns to interesting results about harmonic numbers. Finally, a special family of combinatorial sums is considered and studied in details. To conclude the paper, a few examples of applications to Fibonacci numbers are offered.\\[0.3cm]
For the convenience of the reader, we recall the following definition: 
The harmonic numbers $H_z$ are defined by $H_0=0$ and for $0\ne z\in\mathbb C\setminus\mathbb Z^{-}$ by the relation
\begin{equation*}
H_z = H_{z - 1} + \frac{1}{z}
\end{equation*}
It is worth noticing that
\begin{equation}\label{Har_psi}
H_z = \psi(z + 1) + \gamma,\qquad (z\notin \mathbb{Z}^-)
\end{equation}
where $\gamma$ is the Euler-Mascheroni constant and $\psi(z)=\Gamma'(z)/\Gamma(z)$ is the digamma (or psi) function. 
Of course, when $z=n$, a positive integer, then $H_n=\sum_{k=1}^n \frac{1}{k}$.


\section{A fundamental lemma and first consequences}

The results of this study are derived from the following fundamental lemma.

\begin{lemma}
Let $n$ be a nonnegative integer and let $x$ and $z$ be two complex numbers such that $z\notin\{0,1,\ldots,n\}$. Then we have
\begin{equation}\label{main_lem_id}
\sum_{k=0}^n \binom{z}{k} (-1)^k (x+1)^k = (-1)^n (n+1)\binom{z}{n+1} \sum_{k=0}^n \binom{n}{k} \frac{x^k}{z-k}.
\end{equation}
\end{lemma}
\begin{proof}
We use the following known relations where $z$ is a complex number 
\[\binom{z}{s}\binom{s}{t}=\binom{z}{t}\binom{z-t}{s-t}\quad (s,t \ \mbox{integers and } \ s\ge t\ge 0)\]
and
\[\sum_{k=0}^p(-1)^k\binom{z}{k}=(-1)^p\binom{z-1}{p}\quad (p\ \mbox{integer and }\ p\ge 0).\]
We calculate
\begin{align*}
\sum_{k=0}^n \binom{z}{k} (-1)^k (x+1)^k &= \sum_{k=0}^n \binom{z}{k} (-1)^k \sum_{j=0}^k \binom{k}{j} x^j \\
&= \sum_{j=0}^n \sum_{k=j}^n \binom{z}{j} \binom{z-j}{k-j} (-1)^k x^j \\
&= \sum_{j=0}^n \binom{z}{j} x^j \sum_{r=0}^{n-j} \binom{z-j}{r} (-1)^{r+j} \\
&= \sum_{j=0}^n \binom{z}{j} (-1)^j x^j (-1)^{n-j} \binom{z-j-1}{n-j} \\
&= (-1)^n \sum_{j=0}^n \binom{z}{j} \binom{z-j-1}{n-j} x^j
\end{align*}
and \eqref{main_lem_id} follows because
\[\binom{z}{j} \binom{z-j-1}{n-j} = \frac{z(z-1)\cdots(z-n)}{z-j}\cdot\frac{1}{j!(n-j)!} = (n+1)\binom{z}{n+1}\binom{n}{j}\frac{1}{z-j}.\]
\end{proof}

Identity \eqref{main_lem_id} is very powerful, as will be shown by the number of interesting combinatorial results obtained in what follows.
As a ``warm up", setting $z=-1$ and $x=-1$ in turn yield the elementary identity
$$\sum_{k=0}^n (x+1)^k = \sum_{k=0}^n \binom{n+1}{k+1} x^k$$
and the well-known decomposition into partial fractions as stated in \eqref{partial} (see also \cite{Boyadzhiev})
\begin{equation}\label{eq_9} 
\frac{ n!}{z(z-1)\cdots (z-n)} = \sum_{k=0}^n \binom{n}{k} \frac{(-1)^{n-k}}{z-k}.
\end{equation}
In a way, identity \eqref{main_lem_id} can be considered as generalizing either of these results.\\[0.3cm]

In general, from \eqref{main_lem_id} (or a modified version), we will deduce identities with only one complex parameter and obtain closed forms of combinatorial sums (or identities about them). To illustrate this scheme, here is an introductory example.

\begin{proposition}
Let $n$ be a nonnegative integer and let $x$ and $z$ be two complex numbers such that 
$z\notin\{-1,-2,\ldots,-(n+1)\}$. Then we have
\begin{equation}\label{main_2}
\sum_{k=0}^n \binom{z+k}{k} (x+1)^k = (n+1)\binom{z+n+1}{n+1} \sum_{k=0}^n \binom{n}{k} \frac{x^k}{z+k+1}.
\end{equation}
\end{proposition}
\begin{proof}
We apply \eqref{main_lem_id} with $-(z+1)$ instead of $z$ and use $\binom{-(z+1)}{k}=(-1)^k\binom{z+k}{k}$.
\end{proof}

The following polynomial identity in the variable $z$ is immediately obtained by taking $x=0$:

\begin{corollary}\label{main_2_cor_1 }
If $z$ is a complex number, then for any integer $n\ge 0$
\begin{equation}
\sum_{k=0}^n \binom{z+k}{k} = \binom{z+n+1}{n}.
\end{equation}
In particular, for any nonnegative integer $q$, we have
\begin{equation*}
\sum_{k=0}^n \binom{q+k}{k} = \binom{q+n+1}{n},
\end{equation*}
the well-known hockey-stick formula (see \cite{Lynch} for a different generalization).
\end{corollary}

The next corollary evaluates a combinatorial sum.

\begin{corollary} 
If $n$ is a nonnegative integer, then the following identity holds:
\[\sum_{k=0}^n \binom{n}{k} \frac{(-1)^k}{2^k(n+k+1)} = \frac{2^n}{(2n+1)\binom{2n}{n}}.\]
\end{corollary}
\begin{proof}
In \eqref{main_2}, take $z=n,x=-\frac{1}{2}$ and recall the well-known $\sum_{k=0}^n \binom{n+k}{k}\frac{1}{2^k}=2^n$.
\end{proof}

\section{First applications}

We now apply the scheme described and illustrated above to obtain more elaborate identities from \eqref{main_lem_id}. 
Each proposition of this section is the starting point of a new scheme.\\

Our first example is directly obtained by replacing $z$ by $z+1/2$:
\begin{proposition}
Let $x$ and $z$ be two complex numbers with $z\notin\{-1/2,1/2,\ldots,n-1/2\}$. Then we have
\begin{equation}\label{id_C1}
\sum_{k=0}^n \binom{z+1/2}{k} (-1)^k (1+x)^k= (-1)^n 2(n+1) \binom{z+1/2}{n+1}\sum_{k=0}^n \binom{n}{k} \frac{x^k}{2z-2k+1}.
\end{equation}
\end{proposition}

As a corollary we get two polynomial identities.

\begin{corollary}\label{polynomial1}
If $x$ is a complex number, then for any integer $n\ge 0$ we have
\begin{equation}\label{id 1}
\sum_{k=0}^n \binom{2k}{k} 2^{-2k} x^k = 2^{-2n} (2n+1) \binom{2n}{n} \sum_{k=0}^n \binom{n}{k} (-1)^k \frac{(1-x)^k}{2k+1}
\end{equation}
and
\begin{equation}\label{cor_C11}
\sum_{k=0}^n \binom{2k}{k} 2^{-2k} \frac{(1+x)^k}{2k-1} = 2^{-2n} \binom{2n}{n} \sum_{k=0}^n \binom{n}{k} \frac{x^k}{2k-1}.
\end{equation}
\end{corollary}
\begin{proof}
To obtain \eqref{id 1}, we take $z=-1$ in \eqref{id_C1}, use 
$$\binom{-1/2}{k}=(-1)^k 2^{-2k}\binom{2k}{k}$$ 
and then replace $x$ by $x-1$. \\
Second, we take $z=0$ in \eqref{id_C1} and use 
$$\binom{1/2}{k}=(-1)^{k-1}2^{-2k}\frac{1}{2k-1}\binom{2k}{k}.$$ 
In both cases, we simplify with the help of
$$\binom{2(n+1)}{n+1} = \frac{2(2n+1)}{n+1} \binom{2n}{n}.$$
\end{proof}

\begin{remark}
Identity \eqref{id 1} is known as Carlitz's polynomial relation \cite{Carlitz}. Identity \eqref{cor_C11} seems to be new.
\end{remark}

Two known examples from the previous corollary are
$$\sum_{k=0}^n \binom{2k}{k} 2^{-2k} \frac{1}{2k-1} = -2^{-2n} \binom{2n}{n},$$
and
$$\sum_{k=0}^n \binom{n}{k} \frac{(-1)^k}{2k-1} = - \frac{2^{2n}}{\binom{2n}{n}}.$$
The first example appears in Riordan's book \cite[p. 130]{Riordan} and also as equation (31) in \cite{Adegoke}. \\

In a similar way, we obtain

\begin{proposition}
Let $x$ and $z$ be two complex numbers with $z\notin\{1/2,3/2,\ldots,n+1/2\}$. Then we have
\begin{equation}\label{id_C2}
\sum_{k=0}^n \binom{z-1/2}{k}(-1)^k(x+1)^k= (-1)^n 2(n+1) \binom{z-1/2}{n+1}\sum_{k=0}^n \binom{n}{k} \frac{x^k}{2z-2k-1}.
\end{equation}
\end{proposition}

\begin{corollary}\label{polynomial2}
For each complex $x$ we have
\begin{equation}\label{cor_C21}
\sum_{k=0}^n \frac{\binom{n}{k}}{\binom{2k}{k}} 2^{2k} (-1)^{k+1} (1+x)^{n-k} = \sum_{k=0}^n \binom{n}{k} \frac{x^{n-k}}{2k-1}.
\end{equation}
\end{corollary}
\begin{proof}
Set $z=n$ in \eqref{id_C2} and use the easily checked formulas:
$$\binom{n-1/2}{k}=2^{-2k}\frac{\binom{2n}{n}\binom{n}{k}}{\binom{2n-2k}{n-k}}\qquad (k=0,1,\ldots,n)$$
and
$$\binom{n-1/2}{n+1}=-\frac{2^{-2n-1}}{n+1}\binom{2n}{n}.$$
\end{proof}

As an explicit evaluation we offer the sum
$$\sum_{k=0}^n \frac{\binom{n}{k}}{\binom{2k}{k}} 2^{2k} (-1)^{k+1} = \frac{1}{2n-1}.$$
This sum is also given by equation (32) in \cite{Adegoke}.\\ 

The polynomial identities stated in Corollaries \ref{polynomial1} and \ref{polynomial2} can be merged in two
different ways, hereby leading to additional interesting relations and evaluations.

\begin{corollary}
For each complex $x\neq -1$ we have the following relation
\begin{equation}\label{id_C3}
\sum_{k=0}^n \binom{2k}{k} \frac{2^{-2k}}{2k-1} \left (\frac{1}{1+x}\right )^{n-k}
= 2^{-2n} \binom{2n}{n} \sum_{k=0}^n \frac{\binom{n}{k}}{\binom{2k}{k}} 2^{2k} (-1)^{k+1}\left ( \frac{x}{1+x}\right )^{k}.
\end{equation}
In particular,
\begin{equation*}
\sum_{k=0}^n \binom{2k}{k} \frac{2^{-k}}{2k-1} =2^{-n}\binom{2n}{n}\sum_{k=0}^n\frac{\binom{n}{k}}{\binom{2k}{k}}(-1)^{k+1} 2^{k}
\end{equation*}
and
\begin{equation*}
\sum_{k=0}^n \binom{2k}{k} \frac{2^{-3k}}{2k-1} = - 2^{-3n} \binom{2n}{n} \sum_{k=0}^n \frac{\binom{n}{k}}{\binom{2k}{k}}2^{2k}.
\end{equation*}
\end{corollary}
\begin{proof}
From \eqref{cor_C11} we deduce that
\[ \sum_{k=0}^n \binom{2k}{k}\frac{2^{-2k}}{2k-1}(1+x)^{k-n}=2^{-2n}\binom{2n}{n}(1+x)^{-n}x^n\sum_{k=0}^n\binom{n}{k}\frac{1}{2k-1}\frac{1}{x^{n-k}}.\]
Applying \eqref{cor_C21} (with $\frac{1}{x}$ instead of $x$) gives the conclusion. The special cases follow upon setting first $x=1$ and then $x=-1/2$ in \eqref{id_C3}.
\end{proof}

We advantageously rewrite \eqref{cor_C11} and \eqref{cor_C21} as follows:

\begin{corollary}
\begin{equation}\label{id_Ca}
\sum_{k=0}^n\frac{\binom{n}{k}}{\binom{2k}{k}}2^{2k}(-1)^{k+1} x^k (1+x)^{n-k} = \sum_{k=0}^n\binom{n}{k}\frac{x^k}{2k-1}
=\frac{2^{2n}}{\binom{2n}{n}}\sum_{k=0}^n\binom{2k}{k} 2^{-2k}\frac{(1+x)^{k}}{2k-1}.
\end{equation}
In particular, we get the following identities involving the Catalan numbers $C_n=\frac{1}{n+1}\binom{2n}{n}$
\begin{equation}\label{id_Cb}
\sum_{k=0}^n \frac{2^{2k}}{\binom{2k}{k}} = \frac{1}{3} \left(1+\frac{2^{2n+1}}{C_n} \right)
\end{equation}
and for $n\ge 1$,
\begin{equation}\label{id_Cc}
\sum_{k=1}^n C_k \frac{2^{-2k}}{2k-1} = \frac{1}{3}\left( 1 - 2^{-2n} C_n \right).
\end{equation}
\end{corollary}
\begin{proof}
In \eqref{cor_C21}, replace $x$ by $\frac{1}{x}$ and multiply by $x^n$. This gives the main statement. In \eqref{id_Ca}, we replace $x$ by $-x$ and integrate from $0$ to $1$ to obtain:
\begin{equation*}%\label{id_Cd}
-\sum_{k=0}^n \frac{\binom{n}{k}2^{2k}}{\binom{2k}{k}}\frac{1}{(n+1)\binom{n}{k}}
= \sum_{k=0}^n\binom{n}{k}\frac{(-1)^k}{(k+1)(2k-1)} = \frac{2^{2n}}{\binom{2n}{n}}\sum_{k=0}^n\frac{\binom{2k}{k}}{2^{2k}(k+1)(2k-1)}.
\end{equation*}
From 
$$\frac{1}{(k+1)(2k-1)} = \frac{2}{3}\cdot\frac{1}{2k-1} - \frac{1}{3}\cdot\frac{1}{k+1},$$ 
we then deduce that
\begin{align*}
\sum_{k=0}^n\binom{n}{k}\frac{(-1)^k}{(k+1)(2k-1)} &= \frac{2}{3}\sum_{k=0}^n\binom{n}{k}\frac{(-1)^k}{2k-1}-\frac{1}{3}\sum_{k=0}^n\binom{n}{k}\frac{(-1)^k}{k+1} \\
&= -\frac{2}{3}\frac{2^{2n}}{\binom{2n}{n}}-\frac{1}{3(n+1)}
\end{align*}
(on the right, the first sum has been met earlier and the second sum follows from the obvious $\sum_{k=0}^n\binom{n}{k}\frac{x^{k+1}}{k+1}=\frac{(x+1)^{n+1}-1}{n+1}$).
\end{proof}

\begin{remark}
Identity \eqref{id_Cb} is a classical result that appeared in a range of papers. It was generalized in a different way in \cite{Bataille2}.
\end{remark}

We now modify \eqref{main_lem_id} \textit{via} integration relatively to $x$: 

\begin{proposition}
Let $n$ be a nonnegative integer and let $x$ and $z$ be two complex numbers such that $z\notin\{0,1,\ldots,n\}$. Then we have
\begin{equation}\label{id_Cz}
\sum_{k=0}^n\binom{z}{k}\frac{(-1)^k}{k+1}(1+x)^{k+1}=\sum_{k=0}^n\binom{z}{k}\frac{(-1)^k}{k+1}+(-1)^n(n+1)\binom{z}{n+1}\sum_{k=0}^n\binom{n}{k}\frac{x^{k+1}}{(k+1)(z-k)}
\end{equation}
\end{proposition}
\begin{proof}
In \eqref{main_lem_id}, consider $x$ as a real variable and integrate from $0$ to the  real $X$, to obtain a polynomial identity in $X$; then, rename $X$ as $x$. 
\end{proof}

\begin{corollary}
If $z$ is a complex number with $z\neq -1$, then we have
\begin{equation}\label{id_Cx}
\sum_{k=0}^n \binom{z}{k} \frac{(-1)^k}{k+1} = \frac{1}{z+1}\left (1+(-1)^n \binom{z}{n+1}\right ).
\end{equation}
In addition, if $z\notin\{-1,0,1,\ldots,n\}$, then 
\begin{equation}\label{id_Cy}
\sum_{k=0}^{\lfloor{n/2}\rfloor} \binom{z}{2k} \frac{1}{2k+1} = \frac{1}{z+1} \binom{z}{n+1} \left(\frac{(-1)^{n}+1}{2} + (-1)^n (n+1) 
\sum_{k=0}^n \binom{n}{k} \frac{(-1)^k 2^k}{z-k} \right ).
\end{equation}
\end{corollary}
\begin{proof}
Taking successively $x=-1$ and $x=-2$ in \eqref{id_Cz} gives 
\begin{equation*}
\sum_{k=0}^n\binom{z}{k}\frac{(-1)^k}{k+1} = (-1)^n(n+1)\binom{z}{n+1}\sum_{k=0}^n\binom{n}{k}\frac{(-1)^k}{(k+1)(z-k)}
\end{equation*}
and
\begin{equation*}
\sum_{k=0}^{\lfloor{n/2}\rfloor}\binom{z}{2k}\frac{1}{2k+1} = (-1)^n(n+1)\binom{z}{n+1}\sum_{k=0}^n\binom{n}{k}\frac{(-1)^k 2^k}{(k+1)(z-k)}.
\end{equation*}
We simplify further according to
$$\frac{1}{(z-k)(k+1)} = \frac{1}{z+1}\left (\frac{1}{z-k}+\frac{1}{k+1} \right ).$$
Identity \eqref{id_Cx} follows as we can calculate
\begin{align*}
\sum_{k=0}^n \binom{z}{k} \frac{(-1)^k}{k+1} &= \frac{(-1)^n (n+1)}{z+1}\binom{z}{n+1}\left ( \sum_{k=0}^n \binom{n}{k} \frac{(-1)^k}{z-k} + 
\sum_{k=0}^n \binom{n}{k} \frac{(-1)^k}{k+1} \right ) \\
&= \frac{(-1)^n (n+1)}{z+1}\binom{z}{n+1}\left ( \frac{(-1)^n}{(n+1)\binom{z}{n+1}} + \frac{1}{n+1} \right ),
\end{align*}
where we have also used \eqref{eq_9}. Identity \eqref{id_Cy} is similarly obtained. 
\end{proof}

For later use note that the replacement of $z$ with $-(z+1)$ in \eqref{id_Cx} shows that it can be equivalently written as 
\begin{equation}\label{eq_8}
\sum_{k=0}^n \binom{z+k}{k} \frac{1}{k+1} = \frac{1}{z} \left(\binom{z+n+1}{n+1} - 1 \right)
\end{equation}
for $z\neq 0$.

\begin{corollary}
The following evaluations hold
\begin{equation}\label{Sn1}
\sum_{k=0}^n \binom{n}{k} \frac{(-1)^k}{(k+1)(n-k+1)} = \frac{1+(-1)^n}{(n+1)(n+2)}
\end{equation}
and
\begin{equation}
\sum_{k=0}^{\lfloor{n/2}\rfloor} \binom{n+1}{2k} \frac{1}{2k+1} = \frac{1}{n+2}\left (2^{n+1} + \frac{(-1)^n-1}{2}\right ).
\end{equation}
\end{corollary} 
\begin{proof}
Both sums are particular cases of \eqref{id_Cx} and \eqref{id_Cy} for $z=n+1$. For the second sum we also used the elementary evaluation
$$\sum_{k=0}^n \binom{n}{k} \frac{(-1)^k 2^k}{n+1-k} = (-1)^n \frac{2^{n+1}-1}{n+1}.$$
\end{proof}

The sum \eqref{Sn1} will be encountered gain in Section 5.\\

We proceed by considering special cases of \eqref{id_Cx}.


\begin{corollary}
Let $z$ be a complex number such that $z\neq -3/2$. Then we have
\begin{equation}\label{from_Cx1}
\sum_{k=0}^n \binom{z+1/2}{k}\frac{(-1)^k}{k+1} = \frac{2}{2z+3}\left ( 1 + (-1)^n \binom{z+1/2}{n+1} \right ).
\end{equation}
Also, if $z$ is a complex number such that $z\neq -1/2$, then
\begin{equation}\label{from_Cx2}
\sum_{k=0}^n \binom{z-1/2}{k}\frac{(-1)^k}{k+1}= \frac{2}{2z+1}\left ( 1 + (-1)^n \binom{z-1/2}{n+1} \right ).
\end{equation}
In particular, we have
\begin{equation}
\sum_{k=0}^n 2^{-2k}C_k=2-2^{-2n-1}(2n+1)C_n
\end{equation}
\begin{equation}
\sum_{k=0}^n C_k \frac{2^{-2k}}{2k-1} = - \frac{1}{3}\left ( 2 + 2^{-2n} C_n \right )
\end{equation}
and
\begin{equation}
\sum_{k=0}^n \frac{\binom{n}{k}}{\binom{2(n-k)}{n-k}} 2^{-2k} \frac{(-1)^k}{k+1} = \frac{1}{(n+1)(2n+1)}\left ( \frac{2}{C_n} - (-1)^n 2^{-2n}\right ).
\end{equation}
\end{corollary}

\begin{proof}
The two first identities are obtained by taking $z=-1$ and $z=0$ in \eqref{from_Cx1} and using results already seen in the proof of corollary 3.2. The third identity is derived by setting $z=n$ in \eqref{from_Cx2} and applying the results used in the proof of corollary 3.5.
\end{proof}



\section{Consequences involving harmonic numbers}

This section presents applications to relations involving the harmonic numbers.

\begin{proposition}
Let $z$ be a complex number such that the involved harmonic numbers are defined. Then we have
\begin{equation}\label{Har_from_Cx}
\sum_{k=0}^n \binom{z}{k} \frac{(-1)^k}{k+1} H_{z-k} 
= \frac{H_z}{1+z} + \frac{1}{(1+z)^2} + \frac{(-1)^n}{1+z} \binom{z}{n+1} \left ( H_{z-n-1} + \frac{1}{1+z} \right ).
\end{equation}
\end{proposition}
\begin{proof}
Differentiate \eqref{id_Cx} taking into account that
$$\frac{d}{dz} \binom{z}{k} = \binom{z}{k} (H_z - H_{z-k}).$$
When simplifying use \eqref{id_Cx} again.
\end{proof}

\begin{corollary}
The following evaluations hold
\begin{equation}\label{Hxxx_id1}
\sum_{k=0}^n \binom{n}{k} \frac{(-1)^k}{k+1} H_{n-k} = \frac{1}{n+1}\left ( H_n + \frac{1-(-1)^n}{n+1}\right ),
\end{equation}
\begin{equation}\label{Hxxx_id2}
\sum_{k=0}^n \binom{n}{k} \frac{(-1)^k}{(k+1)(n-k+1)} H_{n+1-k} = \frac{1}{(n+1)(n+2)}\left ( H_{n+1} + \frac{1+(-1)^n}{n+2}\right ),
\end{equation}
and also
\begin{equation}\label{Hxxx_id3}
\sum_{k=0}^n \binom{n}{k} \frac{(-1)^k}{(k+1)(n-k+1)} H_{k+1} = \frac{(-1)^n}{(n+1)(n+2)}\left ( H_{n+1} + \frac{1+(-1)^n}{n+2}\right ).
\end{equation}
\end{corollary}
\begin{proof}
If we take $z=n$ in \eqref{Har_from_Cx}, all terms are defined except $H_{z-n-1}$. However, we have
\[\binom{z}{n+1}H_{z-n-1}=\frac{z-n}{z+1}\binom{z+1}{n+1}H_{z-n-1}=\frac{\binom{z+1}{n+1}}{z+1}\cdot (z-n)H_{z-n-1}\]
and, using \eqref{Har_psi},
\[\lim_{z\to n} (z-n)H_{z-n-1}=\lim_{z\to n} [(z-n)(\psi(z-n)+\gamma)]=-1\]
where the last equality follows from the fact that $0$ is a simple pole of $\psi(z)$ with residue -1. Thus, 
\[\lim_{z\to n}\binom{z}{n+1}H_{z-n-1}=\frac{-1}{n+1}\]
and \eqref{Hxxx_id1} is deduced by a passage to the limit as $z$ tends to $n$ in \eqref{Har_from_Cx}.\\
The second sum is obtained by setting $z=n+1$ in \eqref{Har_from_Cx}. The third sum follows easily form the second.
\end{proof}

The simple polynomial identity that follows is the source of a new series of results.

\begin{proposition}
The following polynomial identity holds
\begin{equation}\label{eq1}
\sum_{k=0}^n \frac{(x+1)^{k+1}}{k+1} = H_{n+1} + \sum_{k=0}^n \binom{n+1}{k+1} \frac{x^{k+1}}{k+1}.
\end{equation}
\end{proposition}
\begin{proof}
Take $z=-1$ in \eqref{id_Cz}.
\end{proof}

\begin{corollary}
If $n$ is a nonnegative integer, then
\begin{equation}\label{eq 5}
\sum_{k=0}^n \binom{n}{k} \frac{(-1)^k}{(k+1)^2} = \frac{H_{n+1}}{n+1} \qquad\mbox{and}\qquad 
\sum_{k=0}^n \binom{n}{k} \frac{(-1)^k 2^k}{(k+1)^2} = \frac{O_{\lfloor{\frac{n}{2}}\rfloor+1}}{n+1},
\end{equation}
where $O_m$ denotes the $m$th odd harmonic number: $O_m=\sum_{k=1}^m\frac{1}{2k-1}$.
\end{corollary}
\begin{proof}
Take successively $x=-1$ and $x=-2$ in \eqref{eq1}.
\end{proof}

The next corollary extends the first of these results.

\begin{corollary}
If $n$ is a nonnegative integer and $q$ a positive integer, then
\begin{equation}\label{eq 3}
\sum_{k=0}^n \binom{n+1}{k+1} \frac{(-1)^k}{(k+1)(k+q+1)} =(n+1)\sum_{k=0}^n \binom{n}{k} \frac{(-1)^k}{(k+1)^2(k+q+1)} = \frac{H_{n+1}}{q}-\frac{1}{q^2}\left(1-\frac{1}{\binom{n+q+1}{q}}\right)
\end{equation}
and
\begin{equation}\label{eq 4}
\sum_{k=0}^n \binom{n+1}{k+1} \frac{(-1)^k}{k+q+1} = (n+1)\sum_{k=0}^n \binom{n}{k} \frac{(-1)^k}{(k+1)(k+q+1)} 
= \frac{1}{q}\left(1-\frac{1}{\binom{n+q+1}{q}}\right).
\end{equation}
\end{corollary}
\begin{proof}
In \eqref{eq1}, replace $x$ by $-x$ and multiply out by $x^{q-1}$ to obtain
\[\sum_{k=0}^n\frac{(1-x)^{k+1}x^{q-1}}{k+1}=H_{n+1}x^{q-1}-\sum_{k=0}^n\binom{n+1}{k+1}\frac{(-1)^k x^{k+q}}{k+1}\]
Integrating from $0$ to $1$ yields
\[\sum_{k=0}^n\frac{1}{k+1}\cdot\frac{(q-1)!(k+1)!}{(q+k+1)!}=\frac{H_{n+1}}{q}- \sum_{k=0}^n\binom{n+1}{k+1}\frac{(-1)^k }{(k+1)(k+q+1)}\]
The left-hand side rewrites as 
\[(q-1)!\sum_{k=0}^n \frac{1}{(k+1)(k+2)\cdots (k+q+1)}=\frac{(q-1)!}{q}\left(\frac{1}{ q!}-\frac{1}{(n+2)\cdots (n+q+1)}\right),\]
where we have used the fact that
\[\frac{q}{(k+1)(k+2)\cdots (k+q+1)}=\frac{1}{(k+1)\cdots (k+q)}-\frac{1}{(k+2)\cdots (k+q+1)}.\]
We deduce that
\[\frac{(q-1)!}{q}\left(\frac{1}{ q!}-\frac{1}{q!\binom{n+q+1}{q}}\right)=\frac{H_{n+1}}{q}- \sum_{k=0}^n\binom{n+1}{k+1}\frac{(-1)^k }{(k+1)(k+q+1)}\]
and \eqref{eq 3} follows. To obtain \eqref{eq 4}, we use the decomposition 
\[\frac{1}{(k+1)(k+q+1)}=\frac{1}{q}\left(\frac{1}{k+1}-\frac{1}{k+q+1}\right)\]
and \eqref{eq 5}.
\end{proof}

\begin{remark}
By reindexing we see that \eqref{eq1} has the equivalent form
\begin{equation*}
\sum_{k=1}^n \frac{(1+x)^k}{k} = H_n + \sum_{k=1}^n \binom{n}{k} \frac{x^k}{k}.
\end{equation*}
This is Equation (22) in \cite{Adegoke}. 
\end{remark}

Identity \eqref{eq1} generalizes as follows.

\begin{proposition}
If $n,q$ are integers with $n\ge 0$ and $q\ge 1$, then
\begin{align}\label{eq2}
\sum_{k=0}^n \frac{(x+1)^{k+q}}{(k+1)(k+2)\cdots (k+q)} &= \sum_{k=0}^n \binom{n+1}{k+1}\frac{x^{k+q}}{(k+1)(k+2)\cdots (k+q)} \nonumber \\
& + H_{n+1}\frac{x^{q-1}}{(q-1)!} + \frac{1}{(q-1)!}\sum_{j=1}^{q-1}\frac{x^{q-1-j}}{j}\binom{q-1}{j}\left(1-\frac{1}{\binom{n+1+j}{j}}\right).
\end{align}
\end{proposition}
\begin{proof}
The proof is by induction on $q$. The case $q=1$ is \eqref{eq1}. The induction step from $q$ to $q+1$ is obtained by first integrating \eqref{eq2} from $0$ to $X$ (and renaming $X$ as $x$ afterward). In the end, all boils down to showing that 
\[\sum_{k=0}^n\frac{1}{(k+1)(k+2)\cdots (k+q+1)}+\frac{1}{(q-1)!}\sum_{j=1}^{q-1}\frac{x^{q-j}}{j(q-j)}\binom{q-1}{j}\left(1-\frac{1}{\binom{n+1+j}{j}}\right)\]
is equal to
\[\frac{1}{q!}\sum_{j=1}^q\frac{x^{q-j}}{j}\binom{q}{j}\left(1-\frac{1}{\binom{n+1+j}{j}}\right).\]
This follows from 
$$\sum_{k=0}^n\frac{1}{(k+1)(k+2)\cdots (k+q+1)} = \frac{1}{q\cdot q!}\left(1-\frac{1}{\binom{n+1+q}{q}}\right)$$ 
(met earlier) and 
$$\frac{q}{j(q-j)} \binom{q-1}{j} = \frac{1}{j}\binom{q}{j},$$ 
which is readily checked.
\end{proof}

\begin{thebibliography}{99}
\bibitem[Adegoke]{Adegoke} K. Adegoke, R. Frontczak and Ch. Hsu, Combinatorial identities of three complex parameters and their basic applications, Open J. Math. Anal. 9 (2) (2025), 66--86.
\bibitem[Bataille2]{Bataille2} M. Bataille and R. Frontczak, New sums mixing harmonic numbers and central binomial coefficients, Fibonacci. Quart. 64 (2026), to appear.
\bibitem[Boyadzhiev]{Boyadzhiev} K. H. Boyadzhiev, Notes on the Binomial Transform, Theory and Table, World Scientific, 2018.
\bibitem[Carlitz]{Carlitz} L. Carlitz, Some identities of Bruckman, Fibonacci Quart. 13 (2) (1975), 121--126.
\bibitem[Lynch]{Lynch} M. Lynch and M. Weselcouch, A generalized Hockey Stick Theorem, J. Integer Seq. 28 (2025), Article 25.7.7.
\bibitem[Riordan]{Riordan} J. Riordan, Combinatorial Identities, R. E. Krieger Publishing Co., 1979.
\end{thebibliography}
\end{document}
